\section{DLA by SAE Feature}
\label{app:feature-resolved-dla}

\Cref{sec:feature-resolved-component-decomposition} uses Sparse Readout Prism
terms inside ordinary direct-logit-attribution accounting as an additive
accounting identity. This appendix states the identity. For a scalar readout target
\(q=\sum_t \alpha_t w_t\), Sparse Readout Prism gives
\[
    q = \sum_i \beta_i(q)d_i + r_q ,
\]
where \(d_i\) is an SAE decoder direction and \(r_q\) is the unembedding
reconstruction error for the selected target. At a hidden state \(h\),
\[
    h^\top q = \sum_i \beta_i(q)h^\top d_i + h^\top r_q .
\]
Under the fixed final-normalization convention used for DLA, decompose
\(h\approx\sum_c h_c\). The DLA term for component \(c\) and SAE feature \(i\)
is
\[
    M_{c,i}(q)=\beta_i(q)\,h_c^\top d_i,
\]
with component reconstruction error term \(R_c(q)=h_c^\top r_q\). We therefore
report two separate gaps: the DLA additivity error between \(h^\top q\) and
\(\sum_c h_c^\top q\), and the readout SAE reconstruction gap between
\(h^\top q\) and \(\sum_{c,i}M_{c,i}(q)\).

The main display uses the pairwise logit difference
\(q=w_{\mathrm{verify}}-w_{\mathrm{assume}}\), for which
\(\beta_i(q)=z_{\mathrm{verify},i}-z_{\mathrm{assume},i}\). The same identity
also covers a row-centered token readout \(q=w_{\mathrm{verify}}-\bar{w}\), which
asks how residual-stream components account for one resolved token's readout
above the vocabulary-mean row. The component view is additive accounting over a
realized forward pass and can be paired with separate interventions for circuit
analysis \citep{geiger2023causalabstraction}.

\Cref{fig:qwen-prism-dla-verify-token-app} shows the row-centered token-readout
version. This is the same source-verification prompt, but the target is
\(h^\top(w_{\mathrm{verify}}-\bar{w})\); the main figure uses the token-token
logit difference.

\begin{figure*}[!tbp]
    \centering
    \AppendixWideGraphic[height=0.34\textheight]{figures/proto_token_lens/qwen2b_32x_prism_dla_verify_token/qwen2b_32x_prism_dla_verify_token_layer_readout.pdf}
    \caption{\textbf{Row-centered token-readout component attribution.}
    The plot decomposes the row-centered \texttt{verify} readout
    \(h^\top(w_{\mathrm{verify}}-\bar{w})\) for the source-verification prompt by
    residual-stream component. Each row stacks readout terms grouped into
    verification/checking features, remaining sparse features, and the
    unembedding reconstruction error term; black points mark ordinary direct logit
    attribution for that residual-stream component. The row-centered readout is
    \(+17.34\), the component DLA sum is \(+17.34\), and the sparse
    SAE feature sum is \(+17.43\), with relative reconstruction error \(0.005\).}
    \label{fig:qwen-prism-dla-verify-token-app}
\end{figure*}

\Cref{fig:qwen-prism-dla-verify-assume-app} is the pairwise companion to the
row-centered view above. It adds the \texttt{assume} side of the contrast, so
verification/checking and assumption-family terms can be read on the same
component axis.

\begin{figure*}[!tbp]
    \centering
    \AppendixWideGraphic[height=0.34\textheight]{figures/proto_token_lens/qwen2b_32x_prism_dla_verify_assume/qwen2b_32x_prism_dla_layer_semantics.pdf}
    \caption{\textbf{Pairwise logit-difference DLA by SAE feature.}
    The plot decomposes the \texttt{verify}-minus-\texttt{assume} logit difference for
    the source-verification prompt by residual-stream component. Each row stacks sparse
    readout terms grouped into verification/checking features, assumption
    features, remaining sparse features, and the unembedding reconstruction error term;
    black points mark ordinary direct logit attribution for that residual-stream
    component. Positive terms support \texttt{verify} over \texttt{assume};
    negative terms support \texttt{assume}. The exact logit difference is \(+4.31\), the
    component DLA sum is \(+4.36\), and the sparse SAE feature sum is
    \(+4.28\), with relative reconstruction error \(0.009\).}
    \label{fig:qwen-prism-dla-verify-assume-app}
\end{figure*}

\FloatBarrier
